\documentclass[11pt]{amsart}
\usepackage{geometry}               
\geometry{letterpaper}                    

\usepackage{graphicx}
\usepackage{tikz}
\usepackage{amssymb}
\usepackage{pdfsync}
\usepackage{hyperref}
\usepackage{verbatim} % this enables the {comment} environment
\usepackage{epstopdf}
\DeclareGraphicsRule{.tif}{png}{.png}{`convert #1 `dirname #1`/`basename #1 .tif`.png}

\newcommand{\pvector}[1]{
  \begin{pmatrix}
    #1
  \end{pmatrix}} 
\newcommand{\ddirac}[1]{
  \,\boldsymbol{\delta}\!\pvector{#1}\!} 
\renewcommand{\d}{\,{\rm d}} 

\def\R{\mathbb{R}}
\def\T{\mathbb{T}}
\def\N{\mathbb{N}}
\def\Z{\mathbb{Z}}
\def\Co{\mathbb{C}}

\def\C{\mathcal{C}}
\def\Q{\mathcal{Q}}
\def\D{\mathcal{D}}
\def\E{\mathcal{E}}
\def\G{\mathcal{G}}
\def\P{\mathbb{P}}
\def\H{\mathcal{H}}
\def\F{\mathcal{F}}
\def\V{\mathcal{V}}
\def\X{\mathcal{X}}
\def\f{\widehat{f}}

\newtheorem{theorem}{Theorem}
\newtheorem{corollary}[theorem]{Corollary}
\newtheorem{definition}[theorem]{Definition}
\newtheorem{remark}[theorem]{Remark}
\newtheorem{proposition}[theorem]{Proposition}
\newtheorem{lemma}[theorem]{Lemma}
\newtheorem{estimate}[theorem]{Estimate}
\newtheorem{claim}[theorem]{Claim}
\newtheorem{question}[theorem]{Question}
\newtheorem{example}[theorem]{Example}
\newtheorem{conclusion}[theorem]{Conclusion}
\newtheorem*{notation}{Notation}


\title{An inverse problem for convolution 
measures with applications}




\author{Diogo Oliveira e Silva}
\address{
 Departamento de Matemática \\
Instituto Superior Técnico \\
Av. Rovisco Pais \\
1049-001 Lisboa, Portugal}
\email{diogo.oliveira.e.silva@tecnico.ulisboa.pt}

\author{João Pedro Ramos}
\address{
 EPFL Lausanne -- Institute of Mathematics MA C3 535 \\
Route Cantonale \\ 
1015 Lausanne, Switzerland}
\email{joao.ramos@epfl.ch}


\date{\today}      
\subjclass[2020]{42B10}
\keywords{Sharp restriction theory, Monge--Ampère equation, stability, inverse problem, convolution measure.} 

\begin{document}
\begin{abstract}
\end{abstract}

\maketitle

\section{To-do list}
\begin{enumerate}
    \item $d\geq 3$ (JPR)
    \item spheres (DOS)
    \item 1st application
    \item $k\geq 3$ for paraboloids
    \item $k\geq 3$ for spheres
    \item 2nd application
    \item NEW (03/2026): Is the map $T$ from OS--Quilodrán (MPCMS 2020) the gradient of a convex function in the sense of optimal transport? Gromov--Knothe's proof of the isoperimetric inequality as explained in Figalli--Maggi--Pratelli (Invent. Math. 2010) was the inspiration for this question, which may lead to stronger and/or simplified proofs of some of our results. Worth exploring!
\end{enumerate}

\section{Introduction}
A classical theorem of Jörgens ($d=2$), Calabi ($d\leq 5$) and Pogorelov ($d\geq 2$) states that any classical convex solution of $\det D^2u=1$ in $\R^d$ must be a quadratic polynomial. As stated in the second paragraph of \cite{CL03}, Caffarelli extended the result for classical solutions to {\it viscosity solutions}. The latter are viscosity subsolutions which are also viscosity supersolutions. We recall the definitions.
\begin{definition}[\cite{Gu16}]
    Let $u\in C^0(\Omega)$ be a convex function and $0\leq f\in C^0(\Omega)$. The function $u$ is a viscosity subsolution (resp.\@ supersolution) of the equation $\det D^2 u=f$ in $\Omega$ if whenever convex $\phi\in C^2(\Omega)$ and $x_0\in\Omega$ are such that $(u-\phi)(x)\leq (u-\phi)(x_0)$ (resp.\@ $\geq$) for all $x$ in a neighborhood of $x_0$, then necessarily 
    \[\det D^2\phi(x_0)\geq f(x_0)\,\,\,(\text{resp.} \leq)\]
\end{definition}
\noindent No generality is lost in assuming that $\phi(x_0)=u(x_0)$ (this follows from considering $\psi=\phi+u(x_0)-\phi(x_0)$, it implies $u(x)\leq\phi(x)$ and motivates the choice of terminology {\it sub}solution). See \cite[Remark 1.3.3.]{Gu16} for the claim that $\phi$ can be further restricted to the class of strictly convex quadratic polynomials.\\

Let $U\subset\R^d$ be open.
Recall that a convex function $f:U\subset\R^d\to\R$ is necessarily continuous.
By Alexandrov's theorem, $f$ has a second derivative almost everywhere (i.e., $f$ has a second-order Taylor expansion at that point with a local error smaller than any quadratic). This is related to Rademacher's theorem,  which ensures that a Lipschitz continuous function $g:U\subset\R^d\to\R$ is differentiable almost everywhere in $U$. The second derivative $D^2f$ is a nonnegative (matrix-valued) Radon measure which is absolutely continuous with respect to $(d-1)$-dimensional Hausdorff measure \cite{Dud77}.

\section{An inverse result}


Our first result is two-dimensional and  phrased for general convex functions $\psi:\R^2\to\R$, without further regularity assumptions.
\begin{proposition}
    Let $\psi:\R^2\to\R$ be strictly convex and let $\sigma$ denote the projection measure on the surface %\footnote{Check distinction between $\psi$ and $|\cdot|^2+\psi$; see \cite{OSQ20} instead of \cite{OSQ18} and check that $C^2$ is not needed for what we export.}
    $\{(y,\psi(y)):y\in\R^2\}$.
    Assume that the function $(\xi,\tau)\mapsto(\sigma\ast\sigma)(\xi,2\psi(\xi)+\tau)$ is a strictly positive constant. Then  $\psi:\R^2\to\R$ is a (positive definite) quadratic form.
\end{proposition}
    
\begin{proof}
Via a change of variables as in the proof of \cite[Prop. 2.1.(c)]{OSQ18} or \cite[Prop. 2.1.(c)]{OSQ20}, we may write
\begin{equation}\label{eq_conv1}
    (\sigma\ast\sigma)(\xi,2\psi(\xi/2)+\tau)=\int_{\R^2}\ddirac{\tau-\Psi(\xi,y)}\,\d y,
\end{equation}
    where $\Psi(\xi,y):=\psi(\xi/2+y)+\psi(\xi/2-y)-2\psi(\xi/2)$ is a nonnegative function in view of the convexity of $\psi$.
    We are working under the assumption that the function $F_\sigma(\xi,\tau):=(\sigma\ast\sigma)(\xi,2\psi(\xi/2)+\tau)$ equals a strictly positive constant, which we denote by $c_0>0$. For every $\xi\in\R^d$ and $t>0$, via \eqref{eq_conv1} and Fubini's theorem we then have that
    \begin{multline}\label{eq_sublevel}
        c_0\cdot t
    =\int_0^t F_\sigma(\xi,\tau)\,\d\tau
    =\int_0^t \int_{\R^2}\ddirac{\tau-\Psi(\xi,y)}\,\d y\,\d\tau\\
    =\int_{\R^2} {\bf 1}_{\{0\leq\Psi(\xi,y)\leq t\}} (y)\,\d y=|\{y\in\R^2: \Psi(\xi,y)\leq t\}|.
    \end{multline}
    It follows that $\psi$ is a viscosity solution of $\det D^2\psi=\pi^2c_0^{-2}$.
    To see why this is the case, let us start by verifying that $\psi$ is a viscosity subsolution. 
    Fix $\xi_0\in\R^2$ and take $\phi\in C^2$ convex such that $\phi(\xi_0)=\psi(\xi_0)$ and $\phi(\xi)\geq\psi(\xi)$ for all $\xi$ in a neighborhood of $\xi_0$. 
    We must check  
    \begin{equation}\label{eq_tocheck}
        \det D^2\phi(\xi_0)\stackrel{?}{\geq} \pi^2c_0^{-2}.
    \end{equation}
    If $t>0$ is small enough, we then have that
    \[c_0\cdot t=|\{y\in\R^2: \Psi(2\xi_0,y)\leq t\}| \geq 
    |\{y\in\R^2:\phi(\xi_0+y)+\phi(\xi_0-y)-2\phi(\xi_0)\leq t\}|,\]
    which implies 
    \begin{equation}\label{eq_sublevel}
        \frac1t |\{y\in\R^2:\phi(\xi_0+y)+\phi(\xi_0-y)-2\phi(\xi_0)\leq t\}|\leq c_0.
    \end{equation}
    We can reverse engineer \eqref{eq_sublevel}, this time for the function $\phi$ instead of $\psi$, and let  $t\to 0^+$ to check, as in the proof of \cite[Prop.\@ 2.1(d)]{OSQ18},
    that
    \[\frac{\pi}{\sqrt{\det D^2\phi(\xi_0)}}\leq c_0,\]
 which is equivalent to \eqref{eq_tocheck}, as desired.
    We now check that $\psi$ is a viscosity supersolution. (Same -- CHECK)

    Thus $\psi$ is a viscosity solution of $\det D^2\psi=c_0$, as claimed. The proposition now follows from \cite[Theorem 1.1]{CL03} (which in  \cite{CL03} is attributed to earlier work of Caffarelli).
\end{proof} 

\subsection{The case $d=1$}
In this case, we have that 
\begin{equation}\label{eq_1dgoal}
    |\{y\in\R:\Psi(\xi,y)\leq t\}|\gtrsim t^{1/2},
\end{equation}
which in light of \eqref{eq_sublevel} leads to a contradiction.
To verify  \eqref{eq_1dgoal}, start by noting that 
\begin{equation}\label{eq_2ndorder}
    \Psi(\xi,y)=\psi(\xi/2+y)+\psi(\xi/2-y)-2\psi(\xi/2)
=\psi''(\xi/2)y^2+o(y^2),
\end{equation}
whenever $\psi$ is twice differentiable.
If there exists $\xi_0\in\R$ such that $\psi''(\xi_0/2)$ exists and is nonzero, then we finish by noting that \eqref{eq_2ndorder} implies that the sublevel set in \eqref{eq_1dgoal} contains a subset of the form $\{y\in\R: y^2\lesssim_{\xi_0} t\}$, which has measure comparable to $t^{1/2}$.
Otherwise, $\psi''=0$ a.e., and we reason as follows:
\begin{align*}
  0\leq \Psi(\xi,y)
   &=\psi(\xi/2+y)-\psi(\xi/2) -(\psi(\xi/2)-\psi(\xi/2-y))\\
   &=\int_{\xi/2}^{\xi/2+y}\psi'(t)\,\d t-\int_{\xi/2-y}^{\xi/2}\psi'(t)\,\d t\\
   &=\int_{\xi/2}^{\xi/2+y} (\psi'(t)-\psi'(t-y))\,\d t\\
   &=\int_{\xi/2}^{\xi/2+y} \int_{t-y}^t \psi''(s)\d s\,\d t\\
   &=\int_{\xi/2-y}^{\xi/2+y} (y-|s-\xi/2|)\psi''(s)\,\d s\\
   &\leq |y| \int_{\xi/2-y}^{\xi/2+y} \psi''(s)\,\d s,
\end{align*}
where the penultimate line follows from Fubini's theorem. We would like to verify that the last integral is $O(|y|)$, which  follows from \cite[Theorem 3.22]{Fo99} as long as the Radon--Nikodym derivative $\frac{\d\mu}{\d m}$ is finite, where $\mu:=\psi''$ and $m$ denotes the one-dimensional Lebesgue measure. But note that Alexandrov's theorem implies  $\mu\ll m$ a.e. (not nec. -- see Remark below, use instead $\psi''=0$ a.e. to conclude from Theorem 3.22), and so $\frac{\d\mu}{\d m}\in L^1(...)$.

\begin{remark}
    An example of a strictly convex function $\psi:\R\to\R$ for which $\psi''=0$ a.e. is as follows. Consider the probability measure
    \[\mu=\sum_{x\in\mathbb Q}\frac{\delta_x}{2^{\sharp x}}\]
    where $\{x_n\}_{n\in\N}$ is a numbering of the rationals $\mathbb Q\subset\R$, and $\sharp x=n$ if $x=x_n$. Let $f(x):=\int_{-\infty}^x\d\mu$ and define 
    \[\psi(x):=\int_0^x f(t)\,\d t.\]
    We have that $\psi'(x)=f(x)>0$ for all $x\in\R$ and $\psi''=\mu$ is a nonnegative measure. So $\psi$ is a strictly increasing convex function. If it were not strictly convex, then $\psi'=f$ would not change on a whole interval, which is a contradiction in view of the density of $\mathbb Q$ in $\R$.
\end{remark}
\subsection{The case $d\geq 3$}
\subsection{On spheres}
 If convolution is constant on the support of the boundary, then it is a sphere.
\section{An application to sharp restriction theory  }


\section*{Acknowledgements}
DOS  is partially partially funded by Fundação para a Ciência e Tecnologia (FCT), Portugal, through  grant No.\@ UID/4459/2025 and  by the Deutsche Forschungsgemeinschaft 
 (DFG, German Research Foundation) under Germany's Excellence Strategy – EXC-2047/1 – 390685813.
JPR ...

\begin{thebibliography}{03}

\bibitem{CL03} L. Caffarelli, Y. Li,
\newblock {\it An extension to a theorem of Jörgens, Calabi, and Pogorelov.}
\newblock Comm. Pure Appl. Math. {\bf 56} (2003), no.~5, 549--583.

\bibitem{Caf90} L. Caffarelli,
\newblock {\it A localization property of viscosity solutions to the Monge--Ampère equation and their strict convexity.}
\newblock Ann. of Math. (2) {\bf 131} (1990), no.~1, 129--134.

\bibitem{Dud77} R. M. Dudley, 
\newblock {\it On second derivatives of convex functions.}
\newblock Math. Scand. {\bf 41} (1977), no.~1, 159--174.

\bibitem{Fo99} G. B. Folland, 
\newblock Real Analysis.
Modern Techniques and Their Applications. Second edition
\newblock Pure Appl. Math. (N. Y.)
Wiley-Intersci. Publ.
John Wiley \& Sons, Inc., New York, 1999.


\bibitem{Gu16} C. Gutiérrez,
\newblock The Monge-Ampère equation.
Second edition. 
\newblock Progr. Nonlinear Differential Equations Appl., {\bf 89}
Birkhäuser/Springer, [Cham], 2016. 

\bibitem{OSQ18} D. Oliveira e Silva, R. Quilodrán, 
\newblock {\it On extremizers for Strichartz estimates for higher order Schrödinger equations.}
\newblock Trans. Amer. Math. Soc. {\bf 370} (2018), no.~10, 6871--6907.

\bibitem{OSQ20} D. Oliveira e Silva, R. Quilodrán, 
\newblock {\it A comparison principle for convolution measures with applications..}
\newblock Math. Proc. Cambridge Philos. Soc.   {\bf 169} (2020), no.~2, 307--322.




\bibitem{St93} E. M. Stein,
\newblock Harmonic Analysis: Real-Variable Methods, Orthogonality, and Oscillatory Integrals.
\newblock Princeton Univ. Press, Princeton, NJ, 1993.
\end{thebibliography}


\end{document}  